% a07_k01steps_b: part of app:a07_k01steps (input from sec/a07_k01steps.tex): the setting of Balaban's step, the exact background move, the elementary lemmas and the difference-kernel bound.

\subsection{The domain-change estimate at a nonzero background}\label{ssec:a07:dom}

\paragraph{The setting.} We are at the step $k\ge1$ of \cite[Sect.~2]{Ba2}, $d=3$, $L=5$, $t\defeq5^k\eps\le\eps_0$, $\eta\defeq5^{-k}$, in scale-$k$ units (unit blocks on $\eta\Z^3$; $\dist$ is the sup-metric \cite[(1.3)]{Ba1}, and $|\cdot|$ counts unit blocks). The couplings are $e_k=et^{1/2}$, $\lambda_k=\lambda t$; $p=p(t)=b_0(1+\ln t^{-1})^{\tilde p}$ is the small-field threshold and $r=r(t)=R(1+\ln t^{-1})^{\varrho}$, $\varrho>1$, $R>R_0$, the layer width \cite[(2.7), p.~558]{Ba2}; $R$ may be enlarged, and $r\ge R$. The large-field regions $\Lambda^{(k)}_0\supset\Lambda^{(k)}_1\supset\dots\supset\Lambda^{(k)}_7$ are those of \cite[(2.8), pp.~558, 566]{Ba2}: $(\Lambda^{(k)}_{i+1})^c$ is the union of the large blocks within distance $r$ of $(\Lambda^{(k)}_i)^c$; $\Lambda^{(k)}_0\subset\Lambda^{(k-1)\prime}_7$, primed sets being the scale-$(k-1)$ sets read at scale $k$; $B^k(S)$ is the set of fine-lattice points of the blocks of $S$. The two propagator domains of \cite[Prop.~2.6]{Ba2} are
\[
\OmA\defeq B^k(\Lambda^{(k-1)\prime}_2)\ \supset\ \OmC\defeq B^k(\Lambda^{(k)}_2),
\]
the inclusion because $\Lambda^{(k)}_2\subset\Lambda^{(k)}_0\subset\Lambda^{(k-1)\prime}_7\subset\Lambda^{(k-1)\prime}_2$; both are unions of large blocks. The background on $B^k(\Lambda^{(k)}_0)$ is $\Bt=(1-\theta_{k+1})A^{(k)}+\theta_{k+1}B^{(k+1),\eta}$ \cite[p.~579]{Ba2}, where $\theta_{k+1}$ equals $1$ on $B^k(\Lambda^{(k)}_2)$ and varies from $1$ to $0$ on a slice of thickness $<M$ around it (\cite[p.~567]{Ba2} at the step $k$, \cite[(2.96)]{Ba2} at the step $k+1$); $\Bt$ is regular on $\OmA$ in the sense of \cite[(2.23)]{Ba1}, with $|\partial\Bt|\le O(p)$ (\cite[(2.60)]{Ba2} for $A^{(k)}$ on $\OmA$, \cite[(2.94)--(2.98)]{Ba2} for $\Bt$ on $B^k(\Lambda^{(k)}_0)$). For a union $\Omega$ of large blocks and a regular field $A$ on it, $G_k(\Omega,A)=(-\Delta^{N}_{A,\Omega}+m^2t^2+a_kP_k(A))^{-1}$ is the scalar propagator of the step on $\Omega$ with Neumann conditions \cite[(2.22)]{Ba1}, \cite[(1.3)]{Ba4}, $P_k(A)$ the projection onto the transported block constants. For a function $F$ of one or two lattice variables on a set $S$, $\norm F_{1,\alpha;S}$ is the natural two-variable extension of Ba{\l}aban's norm \cite[(1.32), p.~420]{Ba3}: the maximum, over the variables, of the sup of $F$, of $|D_AF|$ in each variable, of $|D_xD_{x'}F|$, and of the transported H\"older quotients of exponent $\alpha$ of these derivatives in each variable; for a localisation $h$ ($h=1$ on a cube, smooth, supported in a fixed neighbourhood), $\norm{hF}_{1,\alpha}\le C_h\norm F_{1,\alpha;\supp h}$. We write $\sigma,\sigma'\in\{0,D,HD\}$ for ``no derivative'', ``the covariant derivative'' and ``the H\"older quotient of the covariant derivative'', with $|0|=0$, $|D|=1$, $|HD|=1+\alpha$.

\paragraph{Standing lower bounds.} $R\ge R_0+3$ and
\begin{equation}\label{eq:a07:D7}
D_7\ \defeq\ 5(r-1)-4\ \ge\ \max(R_0+2,\,r),
\end{equation}
where $R_0\ge1$ is the depth constant of \cite[Prop.~2.1]{Ba1}; since $r\ge R$ these are one more lower bound on $R$ or, equivalently, one more cap on $\eps_0$ of Ba{\l}aban's kind.

\paragraph{Printed inputs.} \textup{(P1)} \cite[Prop.~2.1, (2.24)--(2.26)]{Ba1} (which is the theorem \cite[(1.9)--(1.12)]{Ba4}): for a regular $A$ on $\Omega$ and a bounded source $f$, bounds on $G_k(\Omega,A)f$, $DG_k(\Omega,A)f$ and the H\"older quotient of the latter, with decay $c_0e^{-\delta_0\dist(x,\supp f)}$, at points at distance $\ge R_0$ from $\Omega^c$ (2.24)--(2.25), and for $\Omega_2\subset\Omega$ the difference $(G_k(\Omega,A)-G_k(\Omega_2,A))f$ with the additional factor $e^{-\delta_0\dist(\supp f,\Omega_2^c)-\delta_0\dist(\{x,x'\},\Omega_2^c)}$ (2.26), $\delta_0$, $c_0$, $R_0$ independent of $A$, $k$ and $\Omega$; these hold at every scale $j\le k$ after rescaling, by the same statement. \textup{(P2)} \cite[Props.~2.2--2.3]{Ba1} and \cite[(1.13)--(1.20)]{Ba4}: the exponential decay of the kernel of the block-spin operator $\Delta^{(j)}(\Omega,A)$ and of its domain difference $\Delta^{(j)}(\Omega,A)-\Delta^{(j)}(\Omega_0,A)$, the latter with the two-sided factor $e^{-\delta_0(\dist(x,\Omega^{(j)c})+\dist(x',\Omega^{(j)c}))}$ \cite[(2.27)--(2.29)]{Ba1}; the positivity and exponential decay of the restricted block covariances $C^{(j)}_\Lambda(\Omega,A)=\big((\Delta^{(j)}(\Omega,A)+aL^{-2}P_L)|_\Lambda\big)^{-1}$ with the fixed bare $a$ \cite[(2.31)--(2.34)]{Ba1}, \cite[(1.13), (1.15)--(1.16)]{Ba4}; and the two-sided decay of their domain differences $C^{(j)}_\Lambda(\Omega,A)-C^{(j)}_\Lambda(\Omega_0,A)$ \cite[(1.19)--(1.20)]{Ba4}; $\Lambda$ a union of big blocks, constants $c_2$, $\theta_0$ free of $j$, $\Omega$, $\Omega_0$, $\Lambda$ and $A$. \textup{(P3)} the Combes--Thomas bound of \cref{lem:a07:CT} below for the fixed-spacing resolvent, with no restriction near the boundary; the positivity it needs is \cite[(1.8)]{Ba4}. Below, $\theta_0$ denotes the smaller of the rates of \textup{(P2)} and \textup{(P3)}.

\begin{lemma}[a Combes--Thomas bound at fixed spacing]\label{lem:a07:CT}
Let $\Omega$ be a subset of a lattice of spacing $\eta\le1$ with the sup-metric, and on $\ell^2(\Omega)$ (weights $\eta^d$) let $T=D_A^*D_A+V+\sum_Bc_B\Pi_B$, where $D_A$ is the covariant difference on the bonds of $\Omega$, $V\ge0$ is a multiplication operator, the $\Pi_B$ are orthogonal projections whose ranges are supported in pairwise disjoint sets $B$ of diameter $\le\ell_*$, $\eta\le\ell_*$, $0\le c_B\le\bar a$, and $T\ge\gamma>0$. Put $\theta_0\defeq\min\{\ell_*^{-1},\gamma^{1/2}/(36\sqrt d),\gamma/(12e^2\bar a\ell_*)\}$. Then for all $f,f'\in\ell^2(\Omega)$ and all bond functions $g,g'$, with $\delta$ the distance of the supports (a bond located at its initial point),
\begin{align*}
|\ip f{T^{-1}f'}|&\le2\gamma^{-1}e^{-\theta_0\delta}\norm f\norm{f'},\\
|\ip g{D_AT^{-1}f'}|,\ |\ip f{T^{-1}D_A^*g'}|&\le(2\gamma^{-1/2}+6\sqrt d\,\theta_0\gamma^{-1})\,e^{\theta_0}e^{-\theta_0\delta}\,\norm\cdot\norm\cdot,\\
|\ip g{D_AT^{-1}D_A^*g'}|&\le4e^{\theta_0}e^{-\theta_0\delta}\norm g\norm{g'},
\end{align*}
with no restriction on the position of the supports relative to $\partial\Omega$ (the factor $e^{\theta_0}$ covers locating a bond at either end). In particular $\norm{\mathbf 1_YT^{-1}f}\le2\gamma^{-1}e^{-\theta_0\dist(Y,\supp f)}\norm f$ and $\norm{\mathbf 1_YD_AT^{-1}f}\le(2\gamma^{-1/2}+6\sqrt d\,\theta_0\gamma^{-1})e^{\theta_0}e^{-\theta_0\dist(Y,\supp f)}\norm f$ for every set $Y$. For the operators of \cref{lem:a07:K3} in scale-$1$ units ($\eta=1/L$, unit blocks, $\ell_*=1$, $c_B=a$, $V=m^2t^2l_1^2$) the lower bound $T\ge\gamma$ is \cite[(1.8)]{Ba4}, with $\gamma$ free of $\eta$, $\Omega$ and $A$.
\end{lemma}
\begin{proof}
\emph{Weight.} Let $\rho\defeq\dist(\cdot,\supp f')$ on the sites, $\rho(b)\defeq\rho(b_-)$ on the bonds, and $T_\theta\defeq e^{\theta\rho}Te^{-\theta\rho}$; then $T^{-1}=e^{-\theta\rho}T_\theta^{-1}e^{\theta\rho}$ and $e^{\theta\rho}f'=f'$. \emph{Hopping.} The ends of a bond are at distance $\eta$ and $\rho$ is $1$-Lipschitz, so for $\theta\eta\le1$, $e^{\theta\rho}D_Ae^{-\theta\rho}=D_A+\theta B_1$ with $(B_1u)(b)=\theta^{-1}\eta^{-1}\big(e^{\theta(\rho(b_-)-\rho(b_+))}-1\big)U_bu(b_+)$; since $|\theta^{-1}\eta^{-1}(e^{\theta s}-1)|\le e\le3$ for $|s|\le\eta$ and each site is the final point of $d$ bonds, $\norm{B_1}\le3\sqrt d$. Likewise $e^{\theta\rho}D_A^*e^{-\theta\rho}=(D_A-\theta B_1')^*$ with $\norm{B_1'}\le3\sqrt d$. \emph{The hopping change.} $e^{\theta\rho}D_A^*D_Ae^{-\theta\rho}-D_A^*D_A=\theta(D_A^*B_1-B_1'^*D_A)-\theta^2B_1'^*B_1$, so $|\ip\phi{(\cdots)\psi}|\le3\sqrt d\,\theta(\norm{D_A\phi}\norm\psi+\norm\phi\norm{D_A\psi})+9d\theta^2\norm\phi\norm\psi$. \emph{Pins.} Write $\Pi_B=WW^*$ and let $\rho_B$ be the value of $\rho$ at a point of $B$; then $e^{\theta\rho}\Pi_Be^{-\theta\rho}=(e^{\theta(\rho-\rho_B)}W)(e^{-\theta(\rho-\rho_B)}W)^*$, and since $|\rho-\rho_B|\le\ell_*$ on $B$ this differs from $\Pi_B$ by at most $(2e^x+xe^{2x})x\le2e^2x$ in norm, $x\defeq\theta\ell_*\le1$; the $B$ being disjoint, the change of the pin term has norm $\le2e^2\bar a\,\theta\ell_*$. \emph{Perturbation.} Put $K\defeq T^{-1/2}(T_\theta-T)T^{-1/2}$. Since $D_A^*D_A\le T$, $\norm{D_AT^{-1/2}}\le1$, and $\norm{T^{-1/2}}\le\gamma^{-1/2}$; hence $\norm K\le6\sqrt d\,\theta\gamma^{-1/2}+9d\theta^2\gamma^{-1}+2e^2\bar a\theta\ell_*\gamma^{-1}\le\frac16+\frac1{144}+\frac16<\frac12$ for $\theta\le\theta_0$. \emph{Bounded inverse.} $T_\theta=T^{1/2}(1+K)T^{1/2}$, so $T_\theta^{-1}=T^{-1/2}(1+K)^{-1}T^{-1/2}$ with $\norm{T_\theta^{-1}}\le2\gamma^{-1}$, $\norm{D_AT_\theta^{-1}},\norm{T_\theta^{-1}D_A^*}\le2\gamma^{-1/2}$ and $\norm{D_AT_\theta^{-1}D_A^*}\le2$. \emph{Bounds.} $\ip f{T^{-1}f'}=\ip{e^{-\theta\rho}f}{T_\theta^{-1}f'}$ with $\norm{e^{-\theta\rho}f}\le e^{-\theta\delta}\norm f$; $\ip g{D_AT^{-1}f'}=\ip{e^{-\theta\rho}g}{(D_A+\theta B_1)T_\theta^{-1}f'}$ with $\norm{(D_A+\theta B_1)T_\theta^{-1}}\le2\gamma^{-1/2}+6\sqrt d\,\theta\gamma^{-1}$, and the adjoint form follows from the symmetry of $T$; for the third form take $\rho\defeq\dist(\cdot,\supp g')$, so that $\ip g{D_AT^{-1}D_A^*g'}=\ip{e^{-\theta\rho}g}{(D_A+\theta B_1)T_\theta^{-1}(D_A-\theta B_1')^*g'}$ and $\norm{(D_A+\theta B_1)T_\theta^{-1}(D_A-\theta B_1')^*}\le2+12\sqrt d\,\theta\gamma^{-1/2}+18d\theta^2\gamma^{-1}\le4$. Take $\theta=\theta_0$.
\end{proof}

\begin{lemma}[the background move is exact]\label{lem:a07:C51}
On every bond of $\OmC$, $\Bt=B^{(k+1),\eta}$. Hence $G_k(\OmC,\Bt)=G_k(\OmC,B^{(k+1),\eta})$, $\Delta^{(k)}(\OmC,\Bt)=\Delta^{(k)}(\OmC,B^{(k+1),\eta})$ and $C^{(k)}_\Lambda(\OmC,\Bt)=C^{(k)}_\Lambda(\OmC,B^{(k+1),\eta})$ for every $\Lambda\subset\OmC^{(k)}$.
\end{lemma}
\begin{proof}
$\theta_{k+1}=1$ on $\OmC=B^k(\Lambda^{(k)}_2)$, and $\OmC\subset B^k(\Lambda^{(k)}_0)$, so $\Bt=B^{(k+1),\eta}$ there. The three operators depend on the field only through the transporters on the bonds of their domain (the Neumann Laplacian, \cite[(1.3)]{Ba4}) and on the contours $\Gamma^{(k)}_{y,x}$, which lie in $B^k(y)\subset\OmC$ \cite[(1.4)]{Ba4}.
\end{proof}
So the replacement of \cite[Prop.~2.6]{Ba2} is a pure domain change: $G_k(\OmA,\Bt)\to G_k(\OmC,\Bt)$, with $\dG\defeq G_k(\OmA,\Bt)-G_k(\OmC,\Bt)$ on $\OmC\times\OmC$ the $\delta G_k(\Omega,\Omega_2,\Bt)$ of \cite[p.~414]{Ba3}. The terms containing $\tilde B^{(k)}$ were estimated before it (\cite[p.~579]{Ba2}: ``Proposition I.3.1 can also be applied \dots\ the terms containing the configuration $\tilde B^{(k)}$ \dots\ by $O((L^k\eps)^\kappa)|\Lambda^{(k)}_2|$''), as were the terms of order $>\bar n$ and the terms localised in $B^k((\Lambda^{(k)}_7)^c)$; they are not part of the proposition.

\begin{lemma}[three elementary facts]\label{lem:a07:G}
\begin{enumerate}[label=\textup{(\alph*)},nosep,leftmargin=*]
\item \textup{(layer distances)} For $i\ge0$, $j\ge1$: $\dist(\Lambda^{(k)}_{i+j},(\Lambda^{(k)}_i)^c)\ge j(r-1)$; $\dist(\Lambda^{(k-1)\prime}_7,(\Lambda^{(k-1)\prime}_2)^c)\ge r-1$; and $\dist(B^k(S),B^k(S'))\ge\dist(S,S')-1$ on $\eta\Z^3$. In particular every point within distance $2+\tfrac14+\eta<3$ of $B^k(\Lsev)$ lies at depth $\ge D_7$ in $\OmC$.

\item \textup{(smallness)} For $\delta>0$, $n,m\ge0$ and $\kappa>0$ there is $t_0=t_0(\delta,n,m,\kappa,R,\varrho,b_0,\tilde p,\lambda)$ with $p(t)^n\lambda(t)^{-m}e^{-\delta r(t)}\le t^\kappa$ for $t\le t_0$; and for every $\beta>0$, $p(t)^n\le C_{\beta,n}t^{-\beta}$ for $t\le1$.
\item \textup{(tree sums)} Let $d(\{\square\})$ be the length of a shortest tree connecting given cubes of side $\le2$ \cite[p.~420]{Ba3}. For a fixed $\square_1$, $\sum_{\square_2,\dots,\square_n}e^{-\delta d(\{\square_1,\dots,\square_n\})}\le C_T(n,\delta)$, and for a set $S$ of unit cubes $\sum_{\square}e^{-\delta\dist(\square,S)}\le C(\delta)|S|$.
\end{enumerate}
\end{lemma}
\begin{proof}
(a) If $\dist(w,\Lambda_i^c)\le r$, the large block of $w$ lies within $r$ of $\Lambda_i^c$, so the closed $r$-neighbourhood of $\Lambda_i^c$ lies in $\Lambda_{i+1}^c$. Induction on $j$: for $x\in\Lambda_{i+j+1}$ and $z\in\Lambda_i^c$ nearest to $x$, the point $w$ on a sup-metric geodesic from $z$ to $x$ with $\dist(z,w)=\lfloor r\rfloor\ge r-1$ lies in $\Lambda_{i+1}^c$, so $\dist(x,z)\ge j(r-1)+(r-1)$. The second claim is the first at scale $k-1$ ($5(r(t/5)-1)$, divided by $5$, and $r(t/5)\ge r(t)$). The third holds because a block has sup-diameter $\le1$. For the last claim, the first claim with $i=2$, $j=5$ and the third give $\dist(B^k(\Lsev),\OmC^c)\ge5(r-1)-1$, hence depth $>5(r-1)-4=D_7$ at distance $<3$ from $B^k(\Lsev)$. (b) Split $e^{-\delta r}=e^{-\delta r/2}e^{-\delta r/2}$. With $l=\ln t^{-1}\ge0$: $e^{-\delta r/2}=e^{-(\delta R/2)(1+l)^\varrho}\le e^{-(\kappa+m+1)l}=t^{\kappa+m+1}$ once $(1+l)^{\varrho-1}\ge2(\kappa+m+1)/(\delta R)$, that is, for $t\le t_1$; and $b_0^n(1+l)^{n\tilde p}e^{-(\delta R/2)(1+l)^\varrho}$ is bounded in $l$; absorb this bound and $\lambda^{-m}$ by $t\le t_0\le t_1$. The second clause: $(1+l)^{n\tilde p}e^{-\beta l}$ is bounded in $l$. (c) A minimal spanning tree of the cube centres has length $\le2d(\{\square\})+4n$, so $e^{-\delta d}\le e^{2\delta n}\sum_\tau\prod_{uv\in\tau}e^{-(\delta/2)|c_u-c_v|}$ over the $n^{n-2}$ labelled trees; sum the positions leaf by leaf with $\sum_ze^{-(\delta/2)|z|}\le(1+4/\delta)^3$. The second clause is $\sum_{s\in S}\sum_\square e^{-\delta\dist(\square,s)}\le|S|(1+2/\delta)^3$.
\end{proof}

\begin{lemma}[the difference kernel as a two-variable external field]\label{lem:a07:K3}
Let $\Omega_2\subset\Omega$ be unions of large blocks, $A$ regular on $\Omega$, $\dG\defeq G_k(\Omega,A)-G_k(\Omega_2,A)$ on $\Omega_2\times\Omega_2$, and $d(x)\defeq\dist(x,\Omega_2^c)$. Let $D\ge R_0+2$, and let every evaluation point ($x$, $x'$ and the H\"older partners) have $d\ge D-1$. Then for $\sigma,\sigma'\in\{0,D,HD\}$
\[
\big|\partial^\sigma_x\partial^{\sigma'}_{x'}\dG(x,x')\big|\ \le\ C_{\dG}\,\exp\big[-\delta_{\dG}\big(d(x)+d(x')+|x-x'|\big)\big],
\]
and consequently, for localisations $h,h'$ supported at depth $\ge D$,
\[
\norm{h\,\dG\,h'}_{1,\alpha}\ \le\ C_hC_{\dG}\exp\big[-\delta_{\dG}\big(\dist(\supp h,\Omega_2^c)+\dist(\supp h',\Omega_2^c)+\dist(\supp h,\supp h')\big)\big].
\]
$C_{\dG}$ and $\delta_{\dG}$ depend on the constants of \textup{(P1)--(P3)}, on $L$ and on $\alpha$, and are free of $k$, $\eta$, $\Omega$, $\Omega_2$ and $A$. This is the two-variable, two-sided form of the domain-difference clauses \cite[(2.26)]{Ba1}, \cite[(1.20)]{Ba4}, \cite[(2.5)]{Ba3}.
\end{lemma}

\begin{proof}
\emph{Slices.} For $1\le j\le k+1$ put $G^{[k]}_j(\Omega,A)\defeq(-\Delta^N_{A,\Omega}+m^2t^2+a_jl_j^{-2}P_j(A))^{-1}$, $l_j\defeq5^{j-k}$, $P_j(A)$ the projection onto the transported constants of the scale-$j$ blocks; $G^{[k]}_k=G_k(\Omega,A)$. Rescaling $x=l_j\hat x$, $\hat A=l_j^{1/2}A$ keeps the transporters and gives $G^{[k]}_j(x,x')=l_j^{-1}\hat G_j(\hat x,\hat x')$ with $\hat G_j$ the scale-$j$ propagator in scale-$j$ units, $\hat A$ regular at scale $j$. The block identity at scale $j$ (the $\phi\phi$ block of the Schur inverse of the block operator, with Ba{\l}aban's recursion $a_{j+1}=aa_j/(aL^{-2}+a_j)$ \cite[p.~609]{Ba1}) gives, for $1\le j\le k$,
\[
S_j\defeq G^{[k]}_{j+1}-G^{[k]}_j,\qquad S_j(x,x')=a_j^2l_j^{-1}\sum_{\hat y,\hat y'\in\Omega^{(j)}}\mathsf L_j(\hat x,\hat y)\,C^{(j)}(\hat y,\hat y')\,\mathsf L_j(\hat x',\hat y')^{T},
\]
with the leg $\mathsf L_j(\hat x,\hat y)\defeq(\hat G_j\,Q_j^*(\hat A)\delta_{\hat y})(\hat x)$ and the block covariance $C^{(j)}=(\Delta^{(j)}(\Omega,\hat A)+aL^{-2}P_L)^{-1}$ of (P2). Telescoping, $G_k(\Omega,A)=G^{[k]}_1+\sum_{j=1}^{k-1}S_j$; this is Ba{\l}aban's \cite[(2.6), p.~424]{Ba3} in kernel form.

\emph{Slice bounds.} At points at distance $\ge R_0$ from $\Omega^c$, (P1) at scale $j$ applied to the bounded source $Q_j^*\delta_{\hat y}$ (sup $1$, support $B_j(\hat y)$) gives $|\partial^\sigma\mathsf L_j|\le c_0e^{-\delta_0\dist_j(\hat x,B_j(\hat y))}$ for each $\sigma$ (the depth hypothesis holds since $\dist_j\ge\dist\ge R_0$; the H\"older clause (2.24) has no restriction on the separation), and (P2) gives $|C^{(j)}(\hat y,\hat y')|\le c_2e^{-\theta_0|\hat y-\hat y'|}$. Since $|\hat x-\hat x'|\le\dist(\hat x,B(\hat y))+|\hat y-\hat y'|+\dist(\hat x',B(\hat y'))+2$, keeping half of each exponent for the sums over $\hat y,\hat y'$ (\cref{lem:a07:G}(c)) and rescaling (each $D$ is $l_j^{-1}\hat D$, each H\"older quotient $l_j^{-1-\alpha}$ times the scale-$j$ one) gives
\[
|\partial^\sigma_x\partial^{\sigma'}_{x'}S_j(x,x')|\ \le\ C_Kl_j^{-1-|\sigma|-|\sigma'|}e^{-\delta_K|x-x'|/l_j},\qquad\delta_K\defeq\tfrac12\min(\delta_0,\theta_0),
\]
with $C_K$ free of $k$, $j$, $\eta$, $\Omega$, $A$; and $|\partial^\sigma\partial^{\sigma'}G^{[k]}_1|\le C_1(L)l_1^{-1-|\sigma|-|\sigma'|}e^{-\theta_0|x-x'|/(2l_1)}$ by (P3) in scale-$1$ units (spacing $1/L$; point and bond masses of norm $L^{3/2}$; H\"older quotients over separations $\ge1/L$ bounded by $2L^\alpha\sup|D\cdot|$).

\emph{Difference of slices.} Apply the slice identity on $\Omega$ and on $\Omega_2$ with the same $A$: $\dG=\delta G^{[k]}_1+\sum_{j=1}^{k-1}\delta S_j$. By multilinearity each $\delta S_j$ is $a_j^2l_j^{-1}$ times the sum of three terms, $\delta\mathsf L\,C\,\mathsf L^T$, $\mathsf L_2\,\delta C\,\mathsf L^T$, $\mathsf L_2C_2\,\delta\mathsf L^T$, with $\mathsf L=\mathsf L_j(\Omega)$, $\mathsf L_2=\mathsf L_j(\Omega_2)$ (extended by $0$ for $\hat y\notin\Omega^{(j)}_2$), $\delta\mathsf L=\mathsf L-\mathsf L_2$, $C=C^{(j)}(\Omega)$, $C_2=C^{(j)}(\Omega_2)$ (extended by $0$), $\delta C=C-C_2$. Write $d_j\defeq l_j^{-1}d$. \emph{(i) $\delta\mathsf L$.} For $\hat y\in\Omega^{(j)}_2$, (P1)'s clause (2.26) at scale $j$ gives $c_0e^{-\delta_0(\dist(\hat x,B(\hat y))+\dist(B(\hat y),\Omega_2^c)+d_j(\hat x))}$; for $\hat y\notin\Omega^{(j)}_2$, $\delta\mathsf L=\mathsf L(\hat x,\hat y)\le c_0e^{-\delta_0\dist(\hat x,B(\hat y))}$ by (2.25), and $\dist(\hat x,B(\hat y))\ge d_j(\hat x)$ since $B(\hat y)\subset\Omega_2^c$; in all cases $|\partial^\sigma\delta\mathsf L|\le c_0e^{-(\delta_0/2)(\dist(\hat x,B(\hat y))+d_j(\hat x))}$. \emph{(ii) $\delta C$.} For $\hat y,\hat y'\in\Omega^{(j)}_2$ split $\delta C=[C_{\Omega^{(j)}}(\Omega)-C_{\Omega^{(j)}_2}(\Omega)]+[C_{\Omega^{(j)}_2}(\Omega)-C_{\Omega^{(j)}_2}(\Omega_2)]$, both restricted to $\Omega^{(j)}_2$. The first bracket is $-(PXP)^{-1}PXQX^{-1}P$ with $X=\Delta^{(j)}(\Omega)+aL^{-2}P_L$, $P=\mathbf 1_{\Omega^{(j)}_2}$, $Q=1-P$ (Schur identity); $P_L$ has no entries across $\Omega^{(j)}_2$, a union of $L$-blocks, so $PXQ=P\Delta^{(j)}(\Omega)Q$, whose kernel decays exponentially by (P2) (2.27), and both inverses decay at rate $\theta_0$; hence the bracket is $\le Ce^{-\delta_C(|\hat y-\hat y'|+\tilde d(\hat y)+\tilde d(\hat y'))}$ with $\tilde d(\hat y)\defeq\dist(\hat y,(\Omega^{(j)}_2)^c)$, because the path through $Q$ has length $\ge\tilde d(\hat y)+\tilde d(\hat y')$. The second bracket is the domain difference of the restricted block covariance, which has the same form by (P2), \cite[(1.20)]{Ba4}. If $\hat y\in\Omega^{(j)}_2$ and $\hat y'\notin\Omega^{(j)}_2$ (the only other pattern meeting a nonzero $\mathsf L_2$), $\delta C=C(\hat y,\hat y')\le c_2e^{-(\theta_0/2)(|\hat y-\hat y'|+\tilde d(\hat y))}$. \emph{(iii) Assembly.} Proceed as for $S_j$, carrying the decay factors and using $d_j(\hat x')\le d_j(\hat x)+|\hat x-\hat x'|$ and its symmetric form: each term is $\le Ce^{-\delta'(d_j(\hat x)+d_j(\hat x')+|\hat x-\hat x'|)}$, and after rescaling $|\partial^\sigma\partial^{\sigma'}\delta S_j|\le Cl_j^{-1-|\sigma|-|\sigma'|}e^{-\delta'\mathsf S/l_j}$ with $\mathsf S\defeq d(x)+d(x')+|x-x'|$. \emph{(iv) $\delta G^{[k]}_1$.} In scale-$1$ units use the layer identity $H_\Omega=(H_{\Omega_2}\oplus H_{\Omega\setminus\Omega_2})+D_\Sigma^*D_\Sigma$, $\Sigma$ the bonds of $\Omega$ with exactly one end in $\Omega_2$ and $D_\Sigma$ the covariant difference on them (valid because $P_j$ and $m^2$ are block-diagonal and both pieces are unions of blocks), which gives $(G(\Omega)f)|_{\Omega_2}-G(\Omega_2)f=-(G(\Omega)D_\Sigma^*D_\Sigma(G(\Omega_2)f\oplus0))|_{\Omega_2}$; (P3) on $\Omega$ and on $\Omega_2$ bounds the two factors pointwise at the fixed spacing, with no restriction near $\partial\Omega_2$, for each of the four derivative patterns, and the layer weight $L^{-2}$ per bond of $\Sigma$ summed gives $|\partial^\sigma\partial^{\sigma'}\delta\hat G_1|\le C(L)e^{-(\theta_0/2)(d_1(\hat x)+d_1(\hat x')+|\hat x-\hat x'|)}$; rescale with $l_1$. \emph{(v) Sum over $j$.} With $l_j=5^{-n}$, $n=k-j\ge1$, and $\mathsf S\ge2D-2\ge2$: $\sum_n5^{n(3+2\alpha)}e^{-\delta'5^n\mathsf S}\le e^{-\delta'\mathsf S/2}\sum_n5^{n(3+2\alpha)}e^{-\delta'5^n/2}\le Ce^{-\delta'\mathsf S/2}$, since $5^n\mathsf S\ge(5^n+\mathsf S)/2+5^n/2$ for $\mathsf S\ge2$. The localised form follows from the definition of the norm.
\end{proof}

\paragraph{The interaction terms.} The interaction terms of the stage of \cite[p.~579]{Ba2}, after the removals made there (the $R$-vertices, the terms of order $>\bar n$, the terms containing $\tilde B^{(k)}$, the terms localised in $B^k((\Lsev)^c)$), are a finite sum, over label multisets $\cL$ whose boxes lie within $O(1)$ of $B^k(\Lsev)$ (the terms localised in $B^k((\Lsev)^c)$ having been estimated separately on \cite[p.~579]{Ba2}, and \cite[Prop.~2.6]{Ba2} being stated for the terms localised in $B^k(\Lsev)$), over types $T$ and over orders $o\le\bar n$, of per-label full classes $\cC$ (\cref{def:a07:class}): each $\cC$ is the set of all connected graphs of that type, order and label multiset built from the vertices \cite[(1.6)--(1.15)]{Ba3} of the stage, with its Wick and Taylor weights, with the counterterm vertices \cite[(1.7)]{Ba3} carrying their fixed coefficients \cite[(1.23), (1.29)]{Ba3} and labels, \emph{including the whole-graph counterterm members} (Ba{\l}aban's $G'$ with $G_0=G$, \cite[p.~420]{Ba3}: the lone (1.7) vertex with two external legs is a member of the two-scalar-leg classes). This is the connected expansion of the effective action at fixed order and labels; the one-particle-irreducibility of \cite[p.~430]{Ba3} concerns the counterterm-defining classes, whose coefficients are fixed, not the members. The terms may also be finite products of such classes contracted by fixed many-variable external fields \cite[p.~419]{Ba3}. The scalar lines are $G_A\defeq G_k(\OmA,\Bt)$; the external scalar fields are $E_A\defeq a_kG_Af$ with $\norm f_\infty\le Cp\lambda(t)^{-1/4}$, and the other external fields have norms $\le Cp^a\lambda(t)^{-b}$ with $a,b$ bounded by $\bar n$. All vertices lie in $B^k(\Lsev)\subset\OmC$.

